% ===========================================================================
%  Chapitre 8 — Suites numériques
%  PE 2026, composante ALGÈBRE
% ===========================================================================
\bmtchapitre{Suites numériques}
  {Utiliser les suites numériques pour résoudre des problèmes : raisonnement
   par récurrence, convergence, suites du type $u_{n}=f(n)$ et
   $u_{n+1}=g(u_{n})$, suites arithmétiques, géométriques, stationnaires et
   adjacentes.}
  {Algèbre \textbullet\ environ 12 heures}

Une suite est une liste infinie de nombres, numérotée par les entiers. Cette
définition tient en une ligne et ne dit rien de l'intérêt de l'objet : ce qui
compte, c'est qu'une suite décrit un processus qui se répète — une population
qui croît d'année en année, un capital qui produit des intérêts, une valeur
approchée que l'on affine à chaque étape.

Deux outils font la force de ce chapitre. Le raisonnement par récurrence,
d'abord, qui permet de démontrer une propriété pour une infinité d'entiers en
n'écrivant que deux lignes. L'inégalité des accroissements finis, ensuite, qui
transforme une suite définie par $u_{n+1} = f(u_{n})$ en une suite dont on
maîtrise la distance à sa limite. Le problème d'analyse du baccalauréat se
termine presque toujours par cette combinaison.

% ---------------------------------------------------------------------------
\begin{revision}
Les acquis de première et les chapitres d'analyse suffisent.

\begin{enumerate}[label=\textbf{\arabic*.}, leftmargin=*, itemsep=0.35em]
  \item Une suite arithmétique de premier terme $3$ et de raison $2$ :
        donner $u_{5}$ et $u_{n}$.
  \item Une suite géométrique de premier terme $1$ et de raison $\frac12$ :
        donner $v_{4}$ et $v_{n}$.
  \item Calculer $1+2+3+\cdots+100$.
  \item Rappeler l'énoncé de l'inégalité des accroissements finis.
  \item Résoudre $\left(\frac23\right)^{n} \leq 10^{-3}$.
  \item Que vaut $\lim_{n\to+\infty}q^{n}$ selon les valeurs de $q$ ?
\end{enumerate}
\end{revision}

\begin{automatismes}
Sans calculatrice, sans brouillon.

\begin{enumerate}[label=\textbf{\alph*)}, leftmargin=*, itemsep=0.25em]
  \item $u_{n} = 3n-1$ : donner $u_{0}$ et $u_{10}$
  \item $v_{n} = 2^{n}$ : donner $v_{0}$ et $v_{5}$
  \item Somme $1+2+\cdots+n$
  \item Somme $1+q+q^{2}+\cdots+q^{n}$ pour $q \neq 1$
  \item $\lim\left(\frac12\right)^{n}$
  \item $\lim\left(\frac32\right)^{n}$
  \item $\lim (-1)^{n}$
  \item $\lim \frac{1}{n}$
\end{enumerate}
\end{automatismes}

\begin{corrige}[automatismes]
a) $-1$ et $29$ \quad b) $1$ et $32$ \quad c) $\frac{n(n+1)}{2}$ \quad
d) $\frac{1-q^{n+1}}{1-q}$ \quad e) $0$ \quad f) $+\infty$ \quad
g) pas de limite \quad h) $0$.
\end{corrige}

% ===========================================================================
\section{Définition, modes de génération et représentation graphique}

\begin{definition}[suite numérique]
Une \textbf{suite numérique} est une fonction de $\N$ (ou d'une partie de
$\N$) dans $\R$. L'image de l'entier $n$ se note $u_{n}$ et s'appelle le
\textbf{terme de rang} $n$. La suite elle-même se note $(u_{n})$.
\end{definition}

\begin{propriete}[deux modes de génération]
Une suite est donnée soit \textbf{explicitement}, par une formule
$u_{n} = f(n)$ qui calcule directement chaque terme, soit \textbf{par
récurrence}, par un premier terme et une relation $u_{n+1} = g(u_{n})$ qui
fabrique chaque terme à partir du précédent.
\end{propriete}

\begin{remarque}
La différence est pratique avant d'être théorique. Avec une formule
explicite, $u_{1000}$ se calcule immédiatement ; avec une relation de
récurrence, il faut mille étapes. C'est pourquoi tant d'exercices demandent de
passer de la seconde forme à la première, par l'intermédiaire d'une suite
auxiliaire.
\end{remarque}

\begin{visualisation}[la construction en escalier]
\begin{center}
\begin{tikzpicture}
\begin{axis}[
  width = 9.6cm, height = 6.4cm,
  axis lines = middle, xlabel = {$x$},
  xmin = -0.3, xmax = 5.4, ymin = -0.3, ymax = 4.3,
  xtick = {1, 1.75, 2.125}, xticklabels = {$u_{0}$, $u_{1}$, $u_{2}$},
  ytick = \empty,
  axis line style = {bmtGris},
  tick label style = {font=\footnotesize, color=bmtGris}, clip = true,
]
  \addplot[bmtLaterite, line width=1.2pt, domain=0:4, samples=2] {x};
  \addplot[bmtIndigo, line width=1.3pt, domain=0:4, samples=120]
    {0.5*x + 1.25};
  % escalier u0=1 -> u1=1.75 -> u2=2.125...
  \draw[bmtVetiver, line width=0.9pt] (axis cs:1,0) -- (axis cs:1,1.75)
    -- (axis cs:1.75,1.75) -- (axis cs:1.75,2.125)
    -- (axis cs:2.125,2.125) -- (axis cs:2.125,2.3125)
    -- (axis cs:2.3125,2.3125) -- (axis cs:2.3125,2.406);
  \filldraw[bmtVetiver] (axis cs:1,1.75) circle (1.6pt);
  \filldraw[bmtVetiver] (axis cs:1.75,2.125) circle (1.6pt);
  \filldraw[bmtLaterite] (axis cs:2.5,2.5) circle (2pt);
  \node[bmtLaterite, font=\footnotesize, anchor=north west]
    at (axis cs:2.58,2.38) {$\ell$};
  \node[bmtIndigo, font=\footnotesize, anchor=west]
    at (axis cs:4.08,3.25) {$y = g(x)$};
  \node[bmtLaterite, font=\footnotesize, anchor=south east]
    at (axis cs:3.9,3.95) {$y = x$};
\end{axis}
\end{tikzpicture}
\end{center}

Pour construire les termes d'une suite définie par $u_{n+1} = g(u_{n})$, on
part de $u_{0}$ sur l'axe des abscisses, on monte jusqu'à la courbe — on lit
$u_{1}$ en ordonnée —, puis on se déplace horizontalement jusqu'à la droite
$y = x$ pour ramener cette valeur sur l'axe des abscisses. On recommence. La
figure montre aussitôt ce qu'aucun calcul ne montre : la suite converge vers
le point d'intersection de la courbe et de la droite, c'est-à-dire vers la
solution de $g(x) = x$.
\end{visualisation}

% ===========================================================================
\section{Raisonnement par récurrence}

\begin{theoreme}[principe de récurrence]
Soit $P(n)$ une propriété portant sur l'entier $n$, et $n_{0}$ un entier. Si
\begin{enumerate}[label=\textbf{\arabic*.}, leftmargin=*, itemsep=0.1em]
  \item $P(n_{0})$ est vraie ;
  \item pour tout $n \geq n_{0}$, $P(n)$ vraie entraîne $P(n+1)$ vraie ;
\end{enumerate}
alors $P(n)$ est vraie pour tout entier $n \geq n_{0}$.
\end{theoreme}

\begin{methode}{rédiger une récurrence}
Quatre paragraphes, dans cet ordre, et le correcteur les cherche un par un.
\begin{enumerate}[label=\textbf{\arabic*.}, leftmargin=*, itemsep=0.15em]
  \item \emph{Énoncé.} Écrire explicitement la propriété $P(n)$ à démontrer.
  \item \emph{Initialisation.} Vérifier $P(n_{0})$ par le calcul.
  \item \emph{Hérédité.} Supposer $P(n)$ vraie pour un entier $n$ fixé, et
        démontrer $P(n+1)$ en partant de l'hypothèse.
  \item \emph{Conclusion.} Invoquer le principe de récurrence.
\end{enumerate}
\end{methode}

\begin{avertissement}
L'hypothèse de récurrence doit être \emph{utilisée} dans l'hérédité. Une
démonstration qui n'y fait jamais appel n'est pas une récurrence : soit elle
est fausse, soit la propriété se démontrait directement. Autre faute
fréquente : écrire « supposons que $P(n)$ soit vraie pour tout $n$ », ce qui
suppose déjà le résultat.
\end{avertissement}

\begin{exemple}[une récurrence complète]
Montrons que pour tout $n \in \N$, $4^{n}+6n-1$ est divisible par $9$.

\emph{Propriété.} $P(n)$ : « $9$ divise $4^{n}+6n-1$ ».

\emph{Initialisation.} Pour $n = 0$, $4^{0}+0-1 = 0$, divisible par $9$.

\emph{Hérédité.} Supposons $4^{n}+6n-1 = 9k$ avec $k$ entier. Alors
\[
  4^{n+1}+6(n+1)-1 = 4\left(4^{n}+6n-1\right)-18n+9
  = 9\left(4k-2n+1\right),
\]
qui est bien divisible par $9$.

\emph{Conclusion.} La propriété est vraie pour tout entier naturel $n$.
\end{exemple}

\begin{entrainement}[récurrence]
\exo[\niveau{1}] Démontrer que $\displaystyle\sum_{k=1}^{n}k
= \frac{n(n+1)}{2}$.

\exo[\niveau{2}] Démontrer que $\displaystyle\sum_{k=1}^{n}k^{2}
= \frac{n(n+1)(2n+1)}{6}$.

\exo[\niveau{2}] Démontrer que pour tout $n \geq 1$, $2^{n} \geq n+1$.

\exo[\niveau{3}] Soit $(u_{n})$ définie par $u_{0} = 1$ et
$u_{n+1} = \sqrt{2+u_{n}}$. Démontrer que pour tout $n$,
$1 \leq u_{n} \leq 2$.
\end{entrainement}

\begin{corrige}
\textbf{1.} L'égalité est vraie au rang $1$ : $1 = \frac{1\times2}{2}$. Si
elle vaut au rang $n$, alors
$\sum_{k=1}^{n+1}k = \frac{n(n+1)}{2}+(n+1) = \frac{(n+1)(n+2)}{2}$, qui est
la formule au rang $n+1$.

\textbf{2.} Au rang $1$ : $1 = \frac{1\times2\times3}{6}$. Si la formule vaut
au rang $n$,
\[
  \sum_{k=1}^{n+1}k^{2} = \frac{n(n+1)(2n+1)}{6}+(n+1)^{2}
  = \frac{(n+1)\left[n(2n+1)+6(n+1)\right]}{6}
  = \frac{(n+1)(n+2)(2n+3)}{6},
\]
puisque $2n^{2}+7n+6 = (n+2)(2n+3)$.

\textbf{3.} Au rang $1$ : $2 \geq 2$. Si $2^{n} \geq n+1$, alors
$2^{n+1} = 2\times2^{n} \geq 2(n+1) = 2n+2 \geq n+2$, la dernière inégalité
étant vraie dès que $n \geq 0$.

\textbf{4.} Au rang $0$ : $u_{0} = 1$ et $1 \leq 1 \leq 2$. Si
$1 \leq u_{n} \leq 2$, alors $3 \leq 2+u_{n} \leq 4$, donc
$\sqrt3 \leq u_{n+1} \leq 2$. Comme $\sqrt3 > 1$, l'encadrement demandé est
conservé.
\end{corrige}

% ===========================================================================
\section{Suites majorées, minorées, bornées}

\begin{definition}[majorant, minorant]
La suite $(u_{n})$ est \textbf{majorée} s'il existe un réel $M$ tel que
$u_{n} \leq M$ pour tout $n$ ; \textbf{minorée} s'il existe $m$ tel que
$u_{n} \geq m$ pour tout $n$ ; \textbf{bornée} si elle est à la fois majorée
et minorée.
\end{definition}

\begin{propriete}[monotonie]
$(u_{n})$ est \textbf{croissante} si $u_{n+1} \geq u_{n}$ pour tout $n$. Pour
l'établir, on étudie le signe de $u_{n+1}-u_{n}$ ; si tous les termes sont
strictement positifs, on peut aussi comparer $\dfrac{u_{n+1}}{u_{n}}$ à $1$.
\end{propriete}

\begin{entrainement}[monotonie et bornes]
\exo[\niveau{1}] Étudier le sens de variation de $u_{n} = \dfrac{n}{n+1}$ et
montrer qu'elle est bornée.

\exo[\niveau{2}] Étudier le sens de variation de $v_{n} = \dfrac{2^{n}}{n!}$.

\exo[\niveau{2}] Montrer que $w_{n} = \dfrac{\sin n}{n}$ est bornée.

\exo[\niveau{3}] Soit $u_{0} = 2$ et $u_{n+1} = \frac12\left(u_{n}
+\frac{2}{u_{n}}\right)$. Montrer que $u_{n} \geq \sqrt2$ pour tout $n$, puis
que la suite décroît.
\end{entrainement}

\begin{corrige}
\textbf{5.} $u_{n} = 1-\dfrac{1}{n+1}$ : la suite $\left(\frac{1}{n+1}\right)$
décroît, donc $(u_{n})$ croît. De plus $0 \leq u_{n} < 1$ : la suite est
bornée, minorée par $u_{0} = 0$ et majorée par $1$.

\textbf{6.} $\dfrac{v_{n+1}}{v_{n}} = \dfrac{2}{n+1}$. Ce quotient vaut $2$
pour $n = 0$, vaut $1$ pour $n = 1$, et est strictement inférieur à $1$ dès
que $n \geq 2$. La suite croît donc de $v_{0} = 1$ à $v_{1} = 2$, reste égale
à $2$ en $v_{2}$, puis décroît strictement.

\textbf{7.} Pour tout $n \geq 1$, $\abs{\sin n} \leq 1$, donc
$\abs{w_{n}} \leq \dfrac1n \leq 1$ : la suite est bornée, encadrée par $-1$ et
$1$.

\textbf{8.} Pour $a>0$, $\dfrac12\left(a+\dfrac2a\right)-\sqrt2
= \dfrac{a^{2}-2\sqrt2\,a+2}{2a} = \dfrac{\left(a-\sqrt2\right)^{2}}{2a}
\geq 0$ : tout terme obtenu par cette formule est donc supérieur ou égal à
$\sqrt2$, et $u_{0} = 2 \geq \sqrt2$. Ensuite
$u_{n+1}-u_{n} = \dfrac{2-u_{n}^{2}}{2u_{n}} \leq 0$ puisque
$u_{n}^{2} \geq 2$ : la suite décroît.
\end{corrige}

% ===========================================================================
\section{Convergence et divergence d'une suite}

\begin{definition}[limite d'une suite]
La suite $(u_{n})$ \textbf{converge} vers le réel $\ell$ si tout intervalle
ouvert contenant $\ell$ contient tous les termes de la suite à partir d'un
certain rang. On écrit alors $\limn u_{n} = \ell$. Une suite qui ne converge
pas est dite \textbf{divergente}.
\end{definition}

\begin{theoreme}[convergence monotone]
Toute suite croissante et majorée converge. Toute suite décroissante et
minorée converge.
\end{theoreme}

\begin{avertissement}
Ce théorème affirme l'existence de la limite, jamais sa valeur. Une suite
croissante majorée par $10$ ne converge pas nécessairement vers $10$ : elle
converge vers un réel inférieur ou égal à $10$. C'est une erreur classique, et
elle est lourdement sanctionnée.
\end{avertissement}

\begin{theoreme}[limite d'une suite récurrente]
Si $u_{n+1} = g(u_{n})$, si $g$ est continue sur un intervalle contenant tous
les termes, et si $(u_{n})$ converge vers $\ell$, alors $\ell$ vérifie
$g(\ell) = \ell$.
\end{theoreme}

\begin{demonstration}
La suite $(u_{n+1})$ converge vers $\ell$, comme $(u_{n})$ dont elle n'est que
le décalage. Par continuité de $g$, la suite $\left(g(u_{n})\right)$ converge
vers $g(\ell)$. Or ces deux suites sont égales terme à terme : leurs limites
coïncident.
\end{demonstration}

\begin{methode}{encadrer l'écart à la limite}
C'est le schéma du problème d'analyse au baccalauréat, et il se déroule
toujours de la même façon.
\begin{enumerate}[label=\textbf{\arabic*.}, leftmargin=*, itemsep=0.15em]
  \item Montrer que l'équation $g(x) = x$ admet une solution unique $\alpha$
        sur un intervalle $I$ — continuité, stricte monotonie, changement de
        signe.
  \item Montrer par récurrence que tous les termes restent dans $I$.
  \item Majorer $\abs{g'}$ par un réel $k<1$ sur $I$, puis appliquer
        l'inégalité des accroissements finis entre $u_{n}$ et $\alpha$ :
        $\abs{u_{n+1}-\alpha} \leq k\abs{u_{n}-\alpha}$.
  \item En déduire par récurrence
        $\abs{u_{n}-\alpha} \leq k^{n}\abs{u_{0}-\alpha}$, puis conclure à la
        convergence, le majorant tendant vers $0$.
\end{enumerate}
\end{methode}

% ===========================================================================
\section{Suites explicites et suites récurrentes}

\begin{exemple}[la même suite, deux visages]
Soit $u_{0} = 5$ et $u_{n+1} = \dfrac{u_{n}+3}{4}$. Posons
$v_{n} = u_{n}-1$. Alors
\[
  v_{n+1} = u_{n+1}-1 = \frac{u_{n}+3}{4}-1 = \frac{u_{n}-1}{4}
  = \frac14 v_{n}.
\]
La suite $(v_{n})$ est géométrique de raison $\frac14$ et de premier terme
$v_{0} = 4$ : $v_{n} = 4 \times \left(\frac14\right)^{n}$. On revient alors à
la forme explicite
\[
  u_{n} = 1+4\left(\frac14\right)^{n} = 1+4^{1-n}.
\]
Comme $\left|\frac14\right|<1$, la suite converge vers $1$ — ce que
confirme la résolution de $x = \frac{x+3}{4}$.
\end{exemple}

\begin{methode}{deviner la bonne suite auxiliaire}
Le point fixe $\alpha$, solution de $g(x) = x$, est presque toujours la
translation à effectuer : on pose $v_{n} = u_{n}-\alpha$. Si l'énoncé ne
fournit pas la suite auxiliaire, c'est ce calcul-là qu'il faut tenter en
premier.
\end{methode}

% ===========================================================================
\section{Suites arithmétiques et géométriques}

\begin{definition}[suites arithmétiques et géométriques]
$(u_{n})$ est \textbf{arithmétique} de raison $r$ si $u_{n+1} = u_{n}+r$ pour
tout $n$ ; \textbf{géométrique} de raison $q$ si $u_{n+1} = q\,u_{n}$.
\end{definition}

\begin{propriete}[terme général et sommes]
Pour une suite arithmétique de premier terme $u_{0}$ :
\[
  u_{n} = u_{0}+nr, \qquad
  \sum_{k=0}^{n}u_{k} = (n+1)\,\frac{u_{0}+u_{n}}{2}.
\]
Pour une suite géométrique de raison $q \neq 1$ :
\[
  u_{n} = u_{0}\,q^{n}, \qquad
  \sum_{k=0}^{n}u_{k} = u_{0}\,\frac{1-q^{n+1}}{1-q}.
\]
\end{propriete}

\begin{demonstration}
Pour la somme géométrique, on pose $S = u_{0}+u_{0}q+\cdots+u_{0}q^{n}$ et
l'on calcule $S-qS = u_{0}-u_{0}q^{n+1}$ : tous les termes intermédiaires se
détruisent deux à deux. Comme $q \neq 1$, on divise par $1-q$.
\end{demonstration}

\begin{propriete}[limite d'une suite géométrique]
$\limn q^{n} = 0$ si $\abs q<1$ ; $=+\infty$ si $q>1$ ; $=1$ si $q = 1$ ; et
la suite n'a pas de limite si $q \leq -1$.
\end{propriete}

\begin{entrainement}[arithmétiques et géométriques]
\exo[\niveau{1}] $(u_{n})$ est arithmétique, $u_{3} = 7$ et $u_{8} = 22$.
Déterminer $r$, $u_{0}$ et $u_{n}$.

\exo[\niveau{1}] Calculer $1+3+3^{2}+\cdots+3^{10}$.

\exo[\niveau{2}] $(v_{n})$ est géométrique de raison $\frac23$ et
$v_{0} = 81$. À partir de quel rang $v_{n}<10^{-2}$ ?

\exo[\niveau{2}] Une population de $12\,000$ habitants croît de $2{,}5\,\%$
par an. Exprimer $P_{n}$ et déterminer l'année du doublement.

\exo[\niveau{3}] Soit $u_{0} = 2$ et $u_{n+1} = eu_{n}$. Calculer
$S_{n} = u_{0}+\cdots+u_{n}$ et $P_{n} = u_{1}\times\cdots\times u_{n}$.
\end{entrainement}

\begin{corrige}
\textbf{9.} $u_{8}-u_{3} = 5r = 22-7 = 15$, donc $r = 3$. Puis
$u_{3} = u_{0}+3r$ donne $u_{0} = 7-9 = -2$, et $u_{n} = 3n-2$.

\textbf{10.} C'est la somme des onze premiers termes d'une suite géométrique
de raison $3$ et de premier terme $1$ :
$\dfrac{3^{11}-1}{3-1} = \dfrac{177\,147-1}{2} = 88\,573$.

\textbf{11.} $v_{n} = 81\left(\frac23\right)^{n}$, et l'inéquation
$v_{n}<10^{-2}$ équivaut à $\left(\frac32\right)^{n} > 8100$, soit
$n > \dfrac{\ln 8100}{\ln 1{,}5} \approx 22{,}2$. Le premier rang convenable
est $n = 23$ : on vérifie $v_{22} \approx 1{,}08\times10^{-2}$ et
$v_{23} \approx 7{,}2\times10^{-3}$.

\textbf{12.} $P_{n} = 12\,000\times1{,}025^{n}$. Le doublement s'obtient pour
$1{,}025^{n} = 2$, soit $n = \dfrac{\ln2}{\ln 1{,}025} \approx 28{,}1$ : la
population aura doublé au cours de la vingt-neuvième année.

\textbf{13.} La suite est géométrique de raison $e$ et de premier terme $2$,
donc $u_{n} = 2e^{n}$. La somme vaut
$S_{n} = 2\,\dfrac{e^{n+1}-1}{e-1}$, et le produit
$P_{n} = 2^{n}e^{1+2+\cdots+n} = 2^{n}e^{\frac{n(n+1)}{2}}$.
\end{corrige}

% ===========================================================================
\section{Suites arithmético-géométriques}

\begin{definition}[suite arithmético-géométrique]
Une suite est \textbf{arithmético-géométrique} lorsqu'elle vérifie
$u_{n+1} = a\,u_{n}+b$, avec $a \neq 1$.
\end{definition}

\begin{propriete}[réduction]
Soit $\alpha$ le point fixe, solution de $\alpha = a\alpha+b$, c'est-à-dire
$\alpha = \dfrac{b}{1-a}$. Alors la suite $v_{n} = u_{n}-\alpha$ est
géométrique de raison $a$, et
\[
  u_{n} = \alpha + \left(u_{0}-\alpha\right)a^{n}.
\]
\end{propriete}

\begin{demonstration}
$v_{n+1} = u_{n+1}-\alpha = au_{n}+b-(a\alpha+b) = a\left(u_{n}-\alpha\right)
= a\,v_{n}$.
\end{demonstration}

\begin{entrainement}[arithmético-géométriques]
\exo[\niveau{1}] $u_{0} = 0$ et $u_{n+1} = 3u_{n}+4$ : déterminer $u_{n}$.

\exo[\niveau{2}] $u_{0} = 10$ et $u_{n+1} = 0{,}8u_{n}+30$ : déterminer la
limite et l'interpréter.

\exo[\niveau{3}] Un bassin contient $200$ litres. Chaque jour on retire un
cinquième du contenu puis on ajoute $60$ litres. Modéliser, puis déterminer
le contenu à long terme.
\end{entrainement}

\begin{corrige}
\textbf{14.} Le point fixe de $x \mapsto 3x+4$ est $\ell = -2$. En posant
$v_{n} = u_{n}+2$, on obtient $v_{n+1} = 3v_{n}$ avec $v_{0} = 2$, donc
$v_{n} = 2\times3^{n}$ et $u_{n} = 2\times3^{n}-2$.

\textbf{15.} Le point fixe de $x \mapsto 0{,}8x+30$ est $\ell = 150$. En
posant $v_{n} = u_{n}-150$, on obtient $v_{n} = -140\times0{,}8^{n}$, d'où
$u_{n} = 150-140\times0{,}8^{n}$. Comme $0{,}8^{n}\to0$, la suite converge
vers $150$ : quelle que soit la valeur de départ, le procédé se stabilise à
ce niveau.

\textbf{16.} Si $u_{n}$ désigne le contenu au bout de $n$ jours, retirer un
cinquième puis ajouter $60$ donne $u_{n+1} = \frac45u_{n}+60$, avec
$u_{0} = 200$. Le point fixe est $\ell = 300$ ; en posant
$v_{n} = u_{n}-300$, il vient $v_{n} = -100\times\left(\frac45\right)^{n}$,
donc $u_{n} = 300-100\left(\frac45\right)^{n}$. À long terme, le bassin se
stabilise à $300$ litres.
\end{corrige}

% ===========================================================================
\section{Suites stationnaires et suites adjacentes}

\begin{definition}[suite stationnaire, suites adjacentes]
Une suite est \textbf{stationnaire} si elle est constante à partir d'un
certain rang. Deux suites $(u_{n})$ et $(v_{n})$ sont \textbf{adjacentes} si
l'une croît, l'autre décroît, et si $\limn\left(v_{n}-u_{n}\right) = 0$.
\end{definition}

\begin{theoreme}[théorème des suites adjacentes]
Deux suites adjacentes convergent vers la même limite $\ell$, et pour tout
$n$, $u_{n} \leq \ell \leq v_{n}$.
\end{theoreme}

\begin{demonstration}
La suite $(v_{n}-u_{n})$ décroît et tend vers $0$ : elle est donc positive,
d'où $u_{n} \leq v_{n}$ pour tout $n$. Ainsi $(u_{n})$ est croissante et
majorée par $v_{0}$, donc convergente ; de même $(v_{n})$ est décroissante et
minorée par $u_{0}$. Enfin la différence des limites vaut $0$.
\end{demonstration}

\begin{entrainement}[suites adjacentes]
\exo[\niveau{2}] Montrer que $u_{n} = 1-\frac1n$ et $v_{n} = 1+\frac{1}{n^{2}}$
sont adjacentes.

\exo[\niveau{3}] Soit $u_{n} = \displaystyle\sum_{k=0}^{n}\frac{1}{k!}$ et
$v_{n} = u_{n}+\frac{1}{n\times n!}$. Montrer qu'elles sont adjacentes.
\end{entrainement}

\begin{corrige}
\textbf{17.} $(u_{n})$ croît puisque $\frac1n$ décroît ; $(v_{n})$ décroît
puisque $\frac{1}{n^{2}}$ décroît ; enfin
$v_{n}-u_{n} = \frac1n+\frac{1}{n^{2}} \to 0$, et cette différence est
strictement positive. Les deux suites sont donc adjacentes, de limite commune
$1$.

\textbf{18.} $(u_{n})$ croît, chaque terme ajouté étant strictement positif.
Pour $(v_{n})$,
\[
  v_{n+1}-v_{n} = \frac{1}{(n+1)!}+\frac{1}{(n+1)(n+1)!}-\frac{1}{n\times n!}
  = \frac{-1}{n(n+1)(n+1)!} < 0,
\]
après réduction au même dénominateur : la suite décroît. Enfin
$v_{n}-u_{n} = \frac{1}{n\times n!} \to 0$, et cette différence est positive.
Les deux suites sont adjacentes ; leur limite commune est le nombre $e$.
\end{corrige}

% ===========================================================================
\begin{annales}
La suite arrive en fin de problème d'analyse, une fois la fonction étudiée.
Les énoncés qui suivent sont repris de sujets réellement posés ; leur
provenance est indiquée à chaque fois. Lorsqu'un énoncé s'appuie sur une
fonction étudiée en amont du sujet, ses propriétés utiles sont rappelées et
admises.

\exoann{Baccalauréat, Madagascar, série S, session 2021 — problème 2, partie B}
\emph{On admet que la fonction $f$ du sujet vérifie : $f$ est continue sur
$\intervalle{1}{2}$, $f\!\left(\intervalle{1}{2}\right) \subset
\intervalle{1}{2}$, et $\abs{f'(x)} \leq \frac23$ sur cet intervalle. On note
$\alpha$ l'unique solution de $f(x) = x$ dans $\intervalle{1}{2}$.}

Soit $(U_{n})$ définie par $U_{0} = 1$ et $U_{n+1} = f(U_{n})$.
\begin{enumerate}[label=\textbf{\alph*)}, leftmargin=*, itemsep=0.15em]
  \item Démontrer par récurrence que $U_{n} \in \intervalle{1}{2}$ pour tout
        $n$.
  \item Démontrer que $\abs{U_{n+1}-\alpha} \leq \frac23\abs{U_{n}-\alpha}$,
        puis que $\abs{U_{n}-\alpha} \leq \left(\frac23\right)^{n}$.
  \item En déduire la convergence de $(U_{n})$.
  \item Déterminer le plus petit entier $n_{0}$ tel que $U_{n_{0}}$ soit une
        valeur approchée de $\alpha$ à $10^{-3}$ près.
\end{enumerate}

\exoann{Baccalauréat, Madagascar, série C, session 2012 — problème 2, partie A}
\emph{On admet que la fonction $f$ du sujet est continue sur
$I = \intervalle{2}{3}$, que $f(I) \subset I$, et que
$\abs{f'(x)} \leq \frac12$ sur $I$. On note $\alpha$ l'unique solution de
$f(x) = x$ dans $I$.}

Soit $(u_{n})$ définie par $u_{0} = 2$ et $u_{n+1} = f(u_{n})$.
\begin{enumerate}[label=\textbf{\alph*)}, leftmargin=*, itemsep=0.15em]
  \item Démontrer par récurrence que $u_{n} \in I$ pour tout $n$.
  \item Montrer que $\abs{u_{n+1}-\alpha} \leq \frac12\abs{u_{n}-\alpha}$,
        puis que $\abs{u_{n}-\alpha} \leq \left(\frac12\right)^{n}$.
  \item En déduire que $(u_{n})$ converge vers $\alpha$.
  \item Déterminer $n_{0}$ pour que $u_{n_{0}}$ soit une valeur approchée de
        $\alpha$ à $0{,}1$ près.
\end{enumerate}

\exoann{Baccalauréat, Congo-Brazzaville, série C, session 2022 — exercice 3}
\emph{On admet que la fonction $f$ du sujet envoie $\intervalle{2}{2{,}5}$
dans lui-même, que $-0{,}5 \leq f'(x) < 0$ sur cet intervalle, et que $\alpha$
en est l'unique point fixe.}

Soit $(u_{n})$ définie par $u_{0} = 2$ et $u_{n+1} = f(u_{n})$.
\begin{enumerate}[label=\textbf{\alph*)}, leftmargin=*, itemsep=0.15em]
  \item Montrer que si $u_{n} \in \intervalle{2}{2{,}5}$ alors
        $u_{n+1} \in \intervalle{2}{2{,}5}$.
  \item Montrer que $\abs{u_{n+1}-\alpha} \leq \frac12\abs{u_{n}-\alpha}$,
        puis que $\abs{u_{n}-\alpha} \leq
        \left(\frac12\right)^{n}\abs{u_{0}-\alpha} \leq
        \left(\frac12\right)^{n+1}$.
  \item En déduire la convergence, puis trouver $n_{0}$ tel que
        $\abs{u_{n}-\alpha} \leq 10^{-3}$.
\end{enumerate}

\exoann{Baccalauréat, Madagascar, série A, session 2026 — exercice 1}
Soit $(U_{n})$ définie par $U_{0} = 5$ et
$U_{n+1} = \dfrac{U_{n}+3}{4}$.
\begin{enumerate}[label=\textbf{\alph*)}, leftmargin=*, itemsep=0.15em]
  \item On pose $V_{n} = U_{n}-1$. Montrer que $(V_{n})$ est géométrique et
        préciser sa raison.
  \item Exprimer $V_{n}$ puis $U_{n}$ en fonction de $n$, et donner la limite
        de $(U_{n})$.
  \item Calculer $S_{n} = U_{0}+U_{1}+\cdots+U_{n}$.
\end{enumerate}

\exoann{Baccalauréat, Burkina Faso, série D, session 2023 — exercice 2}
Soit $(U_{n})$ définie par $U_{0} = 2$ et
$\ln\!\left(U_{n+1}\right) = 1+\ln\!\left(U_{n}\right)$.
\begin{enumerate}[label=\textbf{\alph*)}, leftmargin=*, itemsep=0.15em]
  \item Montrer que $U_{n+1} = eU_{n}$, puis que $(U_{n})$ est géométrique.
  \item Exprimer $U_{n}$ en fonction de $n$, étudier sa monotonie et sa
        limite.
  \item Calculer $S_{n} = U_{0}+\cdots+U_{n}$, puis
        $P = U_{1}\times\cdots\times U_{n}$.
\end{enumerate}

\exoann{Baccalauréat, Mauritanie, série C, session 2022 — exercice 5}
\emph{On admet que $I_{n} = \int_{0}^{1}x^{n}e^{-x}\dx$ vérifie
$I_{n+1} = (n+1)I_{n}-e^{-1}$ et $I_{0} = 1-\frac1e$.}

On pose $u_{n} = \dfrac{I_{n}}{n!}$.
\begin{enumerate}[label=\textbf{\alph*)}, leftmargin=*, itemsep=0.15em]
  \item Montrer que $u_{n+1}-u_{n} = \dfrac{-e^{-1}}{(n+1)!}$.
  \item En déduire que
        $e\left(1-u_{n}\right) = \displaystyle\sum_{k=0}^{n}\frac{1}{k!}$.
  \item Sachant que $I_{n} \to 0$, en déduire
        $\displaystyle\lim_{n\to+\infty}\sum_{k=0}^{n}\frac{1}{k!}$.
\end{enumerate}

\exoann{Baccalauréat, Madagascar, série S, session 2022 — problème 2, partie B}
Pour $x \in \intervalle{0}{\frac14}$ et $n \in \N^{*}$, on pose
$S_{n}(x) = 1+4x^{2}+\left(4x^{2}\right)^{2}+\cdots+\left(4x^{2}\right)^{n}$.
\begin{enumerate}[label=\textbf{\alph*)}, leftmargin=*, itemsep=0.15em]
  \item Vérifier que
        $S_{n}(x) = \dfrac{1}{1-4x^{2}}-\dfrac{\left(4x^{2}\right)^{n+1}}{1-4x^{2}}$.
  \item En intégrant sur $\intervalle{0}{\frac14}$, montrer que
        \[
          \frac14+\frac{1}{3\times4^{2}}+\frac{1}{5\times4^{3}}+\cdots
          +\frac{1}{(2n+1)\times4^{n+1}} = \frac14\ln3-I_{n},
        \]
        où $I_{n} = \int_{0}^{1/4}\frac{\left(4x^{2}\right)^{n+1}}{1-4x^{2}}\dx$.
  \item Sachant que $I_{n} \to 0$, calculer la limite de la somme.
\end{enumerate}

\exoann{Baccalauréat, Congo-Brazzaville, série C, session 2023 — exercice 3}
\emph{On admet que $f(x) = \frac{6\ln(x+3)}{x+3}$ est décroissante et tend
vers $0$ en $+\infty$.}

On pose $u_{n} = \displaystyle\int_{n}^{n+1}f(x)\dx$.
\begin{enumerate}[label=\textbf{\alph*)}, leftmargin=*, itemsep=0.15em]
  \item Encadrer $f(x)$ sur $\intervalle{n}{n+1}$, puis établir
        $f(n+1) \leq u_{n} \leq f(n)$.
  \item En déduire la limite de $(u_{n})$.
\end{enumerate}

\exoann{Baccalauréat, Côte d'Ivoire, série D, session 2023 — exercice 5}
\emph{On admet que $f(x) = xe^{-x}$ est croissante sur $\intervalle{0}{1}$ et
que $\alpha \in \intervalleo{0}{1}$ vérifie $f(\alpha) = \frac14$.}

Soit $(u_{n})$ définie par $u_{0} = \alpha$ et
$u_{n+1} = u_{n}e^{-u_{n}}$.
\begin{enumerate}[label=\textbf{\alph*)}, leftmargin=*, itemsep=0.15em]
  \item Montrer que tous les termes sont strictement positifs.
  \item Montrer que $(u_{n})$ est décroissante.
  \item En déduire qu'elle converge, puis déterminer sa limite.
\end{enumerate}

\exoann{Baccalauréat, Côte d'Ivoire, série C, session 2024 — exercice 2}
Déterminer, en justifiant, la limite de la suite définie par
$u_{n} = \left(-\dfrac23\right)^{n}$.
\end{annales}

\begin{corrige}[annales]
\textbf{1. a)} $U_{0} = 1 \in \intervalle{1}{2}$ ; si $U_{n} \in
\intervalle{1}{2}$ alors $U_{n+1} = f(U_{n}) \in \intervalle{1}{2}$ puisque
$f$ envoie l'intervalle dans lui-même.
\textbf{b)} L'inégalité des accroissements finis entre $U_{n}$ et $\alpha$
donne $\abs{f(U_{n})-f(\alpha)} \leq \frac23\abs{U_{n}-\alpha}$, soit
$\abs{U_{n+1}-\alpha} \leq \frac23\abs{U_{n}-\alpha}$. Une récurrence donne
$\abs{U_{n}-\alpha} \leq \left(\frac23\right)^{n}\abs{U_{0}-\alpha}$, et
$\abs{U_{0}-\alpha} \leq 1$.
\textbf{c)} $\left(\frac23\right)^{n} \to 0$, donc $U_{n} \to \alpha$.
\textbf{d)} $\left(\frac23\right)^{n} \leq 10^{-3}$ équivaut à
$n \geq \frac{3\ln10}{\ln1{,}5} \approx 17{,}03$ : le plus petit entier est
$n_{0} = 18$.

\textbf{2.} Le raisonnement est identique avec $k = \frac12$ :
$\abs{u_{n}-\alpha} \leq \left(\frac12\right)^{n}$, la suite converge vers
$\alpha$, et $\left(\frac12\right)^{n} \leq 0{,}1$ dès que
$n \geq \frac{\ln10}{\ln2} \approx 3{,}32$, donc $n_{0} = 4$.

\textbf{3.} La stabilité de l'intervalle est donnée par l'énoncé. Comme
$\abs{f'} \leq \frac12$, l'inégalité des accroissements finis donne le
facteur $\frac12$, puis la récurrence
$\abs{u_{n}-\alpha} \leq \left(\frac12\right)^{n}\abs{u_{0}-\alpha}$ ; enfin
$\abs{u_{0}-\alpha} \leq \frac12$ puisque $u_{0}$ et $\alpha$ sont dans un
intervalle de longueur $\frac12$. La convergence suit, et
$\left(\frac12\right)^{n+1} \leq 10^{-3}$ donne $n \geq 9$.

\textbf{4.} $V_{n+1} = U_{n+1}-1 = \frac{U_{n}+3}{4}-1 = \frac{U_{n}-1}{4}
= \frac14V_{n}$ : suite géométrique de raison $\frac14$, de premier terme
$V_{0} = 4$. Donc $V_{n} = 4\left(\frac14\right)^{n}$ et
$U_{n} = 1+4^{1-n}$, qui tend vers $1$. Enfin
$S_{n} = (n+1)+4\dfrac{1-\left(\frac14\right)^{n+1}}{1-\frac14}
= (n+1)+\frac{16}{3}\left(1-4^{-(n+1)}\right)$.

\textbf{5.} $\ln U_{n+1} = \ln(eU_{n})$ donne $U_{n+1} = eU_{n}$ : la suite
est géométrique de raison $e>1$, donc $U_{n} = 2e^{n}$, croissante, de limite
$+\infty$. La somme vaut $S_{n} = 2\dfrac{e^{n+1}-1}{e-1}$. Enfin
$\ln P = \sum_{k=1}^{n}\left(\ln2+k\right) = n\ln2+\frac{n(n+1)}{2}$, donc
$P = 2^{n}e^{n(n+1)/2}$.

\textbf{6.} $u_{n+1}-u_{n} = \frac{(n+1)I_{n}-e^{-1}}{(n+1)!}
-\frac{I_{n}}{n!} = \frac{-e^{-1}}{(n+1)!}$. En sommant de $0$ à $n-1$ et en
utilisant $u_{0} = I_{0} = 1-\frac1e$, on obtient
$u_{n} = 1-\frac1e\sum_{k=0}^{n}\frac{1}{k!}$, d'où l'égalité annoncée. Comme
$I_{n} \to 0$ et $n! \to +\infty$, $u_{n} \to 0$ et la somme tend vers $e$.

\textbf{7.} $S_{n}(x)$ est une somme géométrique de raison $4x^{2} \neq 1$ sur
$\intervalle{0}{\frac14}$, d'où l'expression demandée. En intégrant terme à
terme, $\int_{0}^{1/4}\left(4x^{2}\right)^{k}\dx
= \frac{4^{k}}{4^{2k+1}(2k+1)} = \frac{1}{(2k+1)4^{k+1}}$, ce qui donne la
somme de gauche ; à droite, $\int_{0}^{1/4}\frac{\dx}{1-4x^{2}}
= \frac14\ln3$. La limite de la somme vaut donc $\frac14\ln3$.

\textbf{8.} Sur $\intervalle{n}{n+1}$, la décroissance donne
$f(n+1) \leq f(x) \leq f(n)$ ; l'intervalle étant de longueur $1$,
l'intégration conserve ces bornes. Comme $f(n) \to 0$, l'encadrement donne
$u_{n} \to 0$.

\textbf{9.} Tous les termes sont positifs car $u_{0} = \alpha>0$ et
$xe^{-x}>0$ pour $x>0$. De plus $e^{-u_{n}}<1$, donc $u_{n+1}<u_{n}$ : la
suite décroît et reste minorée par $0$, donc converge vers un réel $\ell$
vérifiant $\ell = \ell e^{-\ell}$, ce qui impose $\ell = 0$.

\textbf{10.} $\abs{-\frac23} = \frac23<1$, donc la suite converge vers $0$ ;
elle est alternée, ses termes changeant de signe à chaque rang.
\end{corrige}

% ===========================================================================
\begin{jecherche}[le remboursement d'un prêt agricole]
Une coopérative emprunte $12\,000\,000$ d'ariary à un taux mensuel de
$1\,\%$. Elle rembourse chaque mois une mensualité constante de $M$ ariary,
versée juste après que les intérêts du mois ont été ajoutés. On note $C_{n}$
le capital restant dû après le $n$-ième versement, et $C_{0} = 12\,000\,000$.

\begin{enumerate}[label=\textbf{\arabic*.}, leftmargin=*, itemsep=0.3em]
  \item Justifier que $C_{n+1} = 1{,}01\,C_{n}-M$.
  \item Déterminer le point fixe de cette relation et l'interpréter : que se
        passe-t-il si $M = 120\,000$ ?
  \item On prend $M = 300\,000$. Exprimer $C_{n}$ en fonction de $n$.
  \item Au bout de combien de mois le prêt est-il soldé ?
  \item Quel est le coût total du crédit ?
  \item La coopérative souhaite solder en $36$ mois exactement. Quelle
        mensualité doit-elle verser ?
\end{enumerate}
\end{jecherche}

\begin{corrige}[le remboursement d'un prêt agricole]
\textbf{1.} Les intérêts du mois valent $1\,\%$ du capital restant, soit
$0{,}01\,C_{n}$ ; on ajoute puis on retire la mensualité :
$C_{n+1} = C_{n}+0{,}01C_{n}-M = 1{,}01C_{n}-M$.

\textbf{2.} Le point fixe vérifie $\alpha = 1{,}01\alpha-M$, soit
$\alpha = 100M$. Pour $M = 120\,000$, $\alpha = 12\,000\,000 = C_{0}$ : le
capital reste constant, la mensualité ne paie que les intérêts et le prêt
n'est jamais remboursé.

\textbf{3.} Pour $M = 300\,000$, $\alpha = 30\,000\,000$ et
$C_{n}-\alpha$ est géométrique de raison $1{,}01$ :
\[
  C_{n} = 30\,000\,000-18\,000\,000\times1{,}01^{n}.
\]

\textbf{4.} $C_{n} \leq 0$ équivaut à
$1{,}01^{n} \geq \frac{30}{18} = \frac53$, soit
$n \geq \frac{\ln(5/3)}{\ln1{,}01} \approx 51{,}3$ : le prêt est soldé au
cours du $52^{\text{e}}$ mois.

\textbf{5.} La coopérative verse $51$ mensualités pleines et un solde partiel.
En arrondissant, le total versé approche $15\,400\,000$ ariary, soit un coût
du crédit d'environ $3\,400\,000$ ariary.

\textbf{6.} Il faut $C_{36} = 0$, c'est-à-dire
$100M-\left(100M-12\,000\,000\right)1{,}01^{36} = 0$, d'où
\[
  M = \frac{12\,000\,000 \times 1{,}01^{36}}{100\left(1{,}01^{36}-1\right)}
  \approx 398\,600 \ \text{ariary}.
\]
\end{corrige}

% ===========================================================================
\begin{jemeteste}
Répondre par vrai ou faux, en justifiant en une ligne.

\begin{enumerate}[label=\textbf{\arabic*.}, leftmargin=*, itemsep=0.28em]
  \item Toute suite croissante tend vers $+\infty$.
  \item Une suite croissante et majorée par $10$ converge vers $10$.
  \item Si $u_{n+1} = g(u_{n})$ et si $(u_{n})$ converge vers $\ell$ avec $g$
        continue, alors $g(\ell) = \ell$.
  \item Une suite bornée converge.
  \item La suite $\left((-1)^{n}\right)$ est bornée et divergente.
  \item Deux suites adjacentes ont la même limite.
\end{enumerate}
\end{jemeteste}

\begin{corrige}[je me teste]
\textbf{1.} Faux — elle peut être majorée, et converger.
\textbf{2.} Faux — elle converge vers un réel inférieur ou égal à $10$.
\textbf{3.} Vrai — c'est le théorème du point fixe.
\textbf{4.} Faux — $\left((-1)^{n}\right)$ est bornée sans converger.
\textbf{5.} Vrai.
\textbf{6.} Vrai — c'est la conclusion du théorème des suites adjacentes.
\end{corrige}

% ===========================================================================
\begin{commeaubac}
\bmtsource{Baccalauréat, série S, session 2021 — problème 2, partie B}
Soit $f$ la fonction du problème, continue sur $\intervalle{1}{2}$, à valeurs
dans $\intervalle{1}{2}$, et telle que $\abs{f'(x)} \leq \frac23$ sur cet
intervalle. On pose $h(x) = f(x)-x$.

\begin{enumerate}[label=\textbf{\arabic*.}, leftmargin=*, itemsep=0.25em]
  \item À l'aide de l'étude de $h$ sur $\intervalle{1}{2}$, montrer que
        l'équation $h(x) = 0$ admet une solution unique
        $\alpha \in \intervalle{1}{2}$. \hfill\textup{(0,75 pt)}
  \item Soit $(U_{n})$ définie par $U_{0} = 1$ et $U_{n+1} = f(U_{n})$.
        \begin{enumerate}[label=\textbf{\alph*)}, leftmargin=1.4em]
          \item Démontrer par récurrence que $U_{n} \in \intervalle{1}{2}$.
                \hfill\textup{(0,5 pt)}
          \item Démontrer que
                $\abs{U_{n+1}-\alpha} \leq \frac23\abs{U_{n}-\alpha}$, puis
                que $\abs{U_{n}-\alpha} \leq \left(\frac23\right)^{n}$.
                \hfill\textup{(0,5 pt $\times$ 2)}
          \item En déduire l'étude de la convergence de $(U_{n})$.
                \hfill\textup{(0,25 pt)}
          \item Déterminer le plus petit entier $n_{0}$ tel que $U_{n_{0}}$
                soit une valeur approchée de $\alpha$ à $10^{-3}$ près.
                \hfill\textup{(0,25 pt)}
        \end{enumerate}
\end{enumerate}
\end{commeaubac}

\begin{corrige}[comme au baccalauréat]
\textbf{1.} $h$ est continue sur $\intervalle{1}{2}$ ; comme $f$ y prend ses
valeurs, $h(1) = f(1)-1 \geq 0$ et $h(2) = f(2)-2 \leq 0$. De plus
$h'(x) = f'(x)-1 \leq \frac23-1<0$ : $h$ est strictement décroissante. Le
théorème des valeurs intermédiaires donne l'existence, la stricte monotonie
l'unicité.

\textbf{2. a)} $U_{0} = 1$ appartient à l'intervalle ; si $U_{n}$ y
appartient, $U_{n+1} = f(U_{n})$ aussi puisque $f$ y prend ses valeurs.
\textbf{b)} L'inégalité des accroissements finis appliquée entre $U_{n}$ et
$\alpha$ donne le facteur $\frac23$ ; la récurrence donne ensuite
$\abs{U_{n}-\alpha} \leq \left(\frac23\right)^{n}\abs{U_{0}-\alpha}$, et
$\abs{U_{0}-\alpha} \leq 1$.
\textbf{c)} Le majorant tend vers $0$ : la suite converge vers $\alpha$.
\textbf{d)} $\left(\frac23\right)^{n} \leq 10^{-3}$ donne
$n \geq \frac{3\ln10}{\ln1{,}5} \approx 17{,}03$, soit $n_{0} = 18$.
\end{corrige}

% ===========================================================================
\begin{concours}
\bmtsource{D'après un concours d'entrée en classes préparatoires}
On pose $u_{n} = \displaystyle\sum_{k=1}^{n}\frac{1}{k}-\ln n$ pour
$n \geq 1$.

\begin{enumerate}[label=\textbf{\arabic*.}, leftmargin=*, itemsep=0.25em]
  \item Montrer que pour tout $k \geq 1$,
        $\dfrac{1}{k+1} \leq \ln(k+1)-\ln k \leq \dfrac{1}{k}$.
  \item En déduire que $(u_{n})$ est décroissante.
  \item Montrer que $(u_{n})$ est minorée par $0$, puis conclure.
\end{enumerate}
\end{concours}

\begin{corrige}[vers le concours]
\textbf{1.} L'inégalité des accroissements finis appliquée à $\ln$ sur
$\intervalle{k}{k+1}$ donne
$\frac{1}{k+1} \leq \ln(k+1)-\ln k \leq \frac1k$, puisque la dérivée
$\frac1x$ y est comprise entre $\frac{1}{k+1}$ et $\frac1k$.

\textbf{2.} $u_{n+1}-u_{n} = \frac{1}{n+1}-\left(\ln(n+1)-\ln n\right)
\leq 0$ d'après l'inégalité de gauche : la suite décroît.

\textbf{3.} En sommant l'inégalité de droite pour $k$ de $1$ à $n-1$, on
obtient $\ln n \leq \sum_{k=1}^{n-1}\frac1k \leq \sum_{k=1}^{n}\frac1k$, donc
$u_{n} \geq 0$. Décroissante et minorée, la suite converge : sa limite est la
constante d'Euler, environ $0{,}577$.
\end{corrige}

% ===========================================================================
\begin{notepourlaclasse}
Douze heures suffisent à condition de ne pas refaire le programme de
première. Les suites arithmétiques et géométriques doivent être révisées en
une séance, pas trois : ce qui est neuf ici, c'est la récurrence et le
traitement de $u_{n+1} = f(u_{n})$.

Sur la récurrence, la faute constante est l'hérédité qui n'utilise pas
l'hypothèse. Faites-la repérer par les élèves eux-mêmes : distribuez deux
copies, l'une correcte, l'autre non, et demandez de trouver la différence.
L'exercice prend dix minutes et vaut mieux qu'un discours.

Le schéma du problème d'analyse — point fixe, stabilité de l'intervalle,
inégalité des accroissements finis, majoration géométrique — mérite d'être
affiché au tableau une fois pour toutes, sous forme de quatre étapes. Les
élèves reconnaissent alors immédiatement le scénario, quelle que soit la
fonction, et il ne leur reste plus qu'à mener les calculs.
\end{notepourlaclasse}

% =========================================================================
\begin{devoir}[2 heures]
Trois exercices indépendants, notés sur $20$. Le devoir est cumulatif : il porte sur ce chapitre et sur ceux qui le précèdent. Les énoncés sont repris de sessions réelles et assemblés ici en épreuve ; leur provenance est indiquée à chaque fois.

\exods{1}{Baccalauréat, Cameroun, séries C et E, session 2024 — exercice 3}{5 points}
Soit $f(x) = e^{-x^{2}}$ sur $\R$. Étudier les variations de $f$ et en déduire
un encadrement de $f(x)$ pour $x$ appartenant à $\intervalle{0}{1}$.

\exods{2}{Baccalauréat, Côte d'Ivoire, série C, session 2024 — exercice 2}{7 points}
Résoudre dans $\R$ l'équation différentielle $y'+2y = 0$, puis déterminer la
solution vérifiant $y(0) = 3$.

\exods{3}{Baccalauréat, Congo-Brazzaville, série C, session 2022 — exercice 3}{8 points}
\emph{On admet que la fonction $f$ du sujet envoie $\intervalle{2}{2{,}5}$
dans lui-même, que $-0{,}5 \leq f'(x) < 0$ sur cet intervalle, et que $\alpha$
en est l'unique point fixe.}

Soit $(u_{n})$ définie par $u_{0} = 2$ et $u_{n+1} = f(u_{n})$.
\begin{enumerate}[label=\textbf{\alph*)}, leftmargin=*, itemsep=0.15em]
  \item Montrer que si $u_{n} \in \intervalle{2}{2{,}5}$ alors
        $u_{n+1} \in \intervalle{2}{2{,}5}$.
  \item Montrer que $\abs{u_{n+1}-\alpha} \leq \frac12\abs{u_{n}-\alpha}$,
        puis que $\abs{u_{n}-\alpha} \leq
        \left(\frac12\right)^{n}\abs{u_{0}-\alpha} \leq
        \left(\frac12\right)^{n+1}$.
  \item En déduire la convergence, puis trouver $n_{0}$ tel que
        $\abs{u_{n}-\alpha} \leq 10^{-3}$.
\end{enumerate}

\end{devoir}

\begin{corrige}[devoir surveillé]
\textbf{Exercice 1.} $f'(x) = -2xe^{-x^{2}}$, du signe de $-x$ : $f$ croît sur
$\left]-\infty\,;0\right]$ et décroît ensuite, avec pour maximum $f(0) = 1$.
Sur $\intervalle{0}{1}$ elle décroît de $1$ à $e^{-1}$, donc
$\frac1e \leq f(x) \leq 1$.

\textbf{Exercice 2.} Les solutions sont les $x\mapsto Ce^{-2x}$ ; la condition
$y(0) = 3$ impose $C = 3$, soit $y(x) = 3e^{-2x}$.

\textbf{Exercice 3.} La stabilité de l'intervalle est donnée par l'énoncé. Comme
$\abs{f'} \leq \frac12$, l'inégalité des accroissements finis donne le
facteur $\frac12$, puis la récurrence
$\abs{u_{n}-\alpha} \leq \left(\frac12\right)^{n}\abs{u_{0}-\alpha}$ ; enfin
$\abs{u_{0}-\alpha} \leq \frac12$ puisque $u_{0}$ et $\alpha$ sont dans un
intervalle de longueur $\frac12$. La convergence suit, et
$\left(\frac12\right)^{n+1} \leq 10^{-3}$ donne $n \geq 9$.

\end{corrige}
